\chapter{Selected EPL Programs}
\label{app:examples}

This appendix includes representative programs from the repository. They are
used by the test harness and by the thesis discussion.

\section{Hello World}

\begin{lstlisting}[language=EPL,caption={examples/hello\_world.epl}]
fn main() -> i32 {
	printf("hello world\n");
	0
}

fn printf(fmt: *i8, ...);
\end{lstlisting}

\section{Fibonacci}

\begin{lstlisting}[language=EPL,caption={examples/fibonacci.epl}]
fn fib(n: u32) -> u32 {
	if n < 2 {
		n
	} else {
		fib(n - 2) + fib(n - 1)
	}
}

fn fib_it(n: u32) -> u32 {
	let a: u32 = 0;
	let b: u32 = 1;
	let i: u32 = 0;
	while i < n {
		let tmp = a;
		a = b;
		b += tmp;
		i += 1;
	}
	a
}

fn main() -> i32 {
	for n in 0u32..10 {
		printf("fib(%i) = %i\n", n, fib(n));
		printf("fib_it(%i) = %i\n", n, fib_it(n));
	}
	0
}

fn printf(fmt: *i8, ...);
\end{lstlisting}

\section{Compile-Time Primes}

\begin{lstlisting}[language=EPL,caption={examples/comptime\_primes.epl}]
fn main() -> i32 {
	let primes = comptime {
		let lut: [i32; 20];
		let ready = 0;
		let next = 3;
		while ready < 20 {
			while !is_prime(next) {
				next += 2;
			}
			lut[ready as u64] = next;
			ready += 1;
			next += 2;
		}
		lut
	};
	for i in 0u64..20 {
		printf("%i ", primes[i]);
	}
	printf("\n");
	0
}

@pure
fn is_prime(x: i32) -> bool {
	if x <= 1 {
		return false;
	}
	if x == 2 {
		return true;
	}
	if x % 2 == 0 {
		return false;
	}
	let i = 3;
	while i <= x / 2 {
		if x % i == 0 {
			return false;
		}
		i += 2;
	}
	true
}

fn printf(fmt: *i8, ...);
\end{lstlisting}

\section{Conway's Game of Life}

\begin{lstlisting}[language=EPL,caption={examples/game\_of\_life.epl}]
struct Board {
	rows: [[bool; 20]; 20]
}

fn main() -> i32 {
	let board: Board;
	clear_board(&board);
	seed_board(&board);

	loop {
		print_board(&board);
		update_board(&board);
		usleep(250000);
	}
}

fn clear_board(board: *Board) {
	for row in 0u64..20 {
		for col in 0u64..20 {
			set_cell(board, row, col, false);
		}
	}
}

fn seed_board(board: *Board) {
	set_cell(board, 8, 10, true);
	set_cell(board, 9, 11, true);
	set_cell(board, 10, 9, true);
	set_cell(board, 10, 10, true);
	set_cell(board, 10, 11, true);
}

fn print_board(board: *Board) {
	printf("\n");
	for row in 0u64..20 {
		for col in 0u64..20 {
			printf(if get_cell(board, row, col) { "#" } else { "." });
		}
		printf("\n");
	}
}

fn update_board(board: *Board) {
	let next: Board;

	for row in 0u64..20 {
		for col in 0u64..20 {
			let alive = get_cell(board, row, col);
			let neighbors = get_alive_neighbors(board, row, col);
			let next_cell = if alive { neighbors == 2 || neighbors == 3 } else { neighbors == 3 };
			set_cell(&next, row, col, next_cell);
		}
	}

	board.* = next;
}

fn get_alive_neighbors(board: *Board, row: u64, col: u64) -> u64 {
	let count = 0u64;

	for r in (row + 20 - 1)..(row + 20 + 2) {
		for c in (col + 20 - 1)..(col + 20 + 2) {
			let r = r % 20;
			let c = c % 20;
			if (r != row || c != col) && get_cell(board, r, c) {
				count += 1;
			}
		}
	}

	count
}

fn get_cell(board: *Board, row: u64, col: u64) -> bool {
	board.*.rows[row][col]
}

fn set_cell(board: *Board, row: u64, col: u64, alive: bool) {
	board.*.rows[row][col] = alive;
}

fn printf(fmt: *i8, ...);
fn usleep(usec: u32) -> i32;
\end{lstlisting}

